\clearpage
\section{Decision Rules}
\label{app:rules}

This appendix restates the decision rules of the frozen plans; nothing in it was chosen after the data were seen.

\paragraph{Primary test of order moderation (adults).} For each representation $r$, item set (8 within-domain items or all 24), and experiment $x = 1, \dots, 12$, stage~1 fits the binomial mixed model
\[
\operatorname{logit} P(y = 1) = \beta_{0x} + \beta_x\, z_d + \textstyle\sum_c \gamma_{cx}\, C_c + \sum_c \eta_{cx}\, z_d C_c + \lambda_{x,\text{level}} + u_{\text{person}},
\]
where $z_d$ is the distance standardized within representation and item set over all experiments, $C_c$ are sum-to-zero contrasts for the four teaching conditions (so $\beta_x$ is the slope averaged over conditions), and the random-effects structure is the richest rung of a fixed ladder, from (1 + $z_d$~$\mid$~person) + (1~$\mid$~image) down to (1~$\mid$~person), that converged without singularity in all 12 experiments for all 10 representations. That rung was (1~$\mid$~person) for both item sets. Stage~2 computes
\[
\hat\Delta_r = \frac{1}{6}\sum_{x \in \mathcal{E}_{31}} \hat\beta_x - \frac{1}{6}\sum_{x \in \mathcal{E}_{13}} \hat\beta_x ,
\]
the mean slope in the six experiments with the one-example trial last minus the mean in the six with it first. The two-sided $p$ value is the share of the $\binom{12}{6} = 924$ assignments of the order labels for which $|\hat\Delta^{*}| \geq |\hat\Delta_r|$ (smallest attainable value $2/924 \approx .0022$). The 95\% and 90\% intervals collect the shifts $\delta_0$ (grid step .001) for which the test applied to $\hat\Delta_r - \delta_0$ (subtracting $\delta_0$ from the slopes of the one-example-last experiments) does not reject at $\alpha = .05$ and .10, respectively. Because order is confounded with blocking, a rejection would not identify order as the cause.

\paragraph{Regions for the primary test.} With 8 items and for each margin $\delta \in \{.05, .10, .20\}$ (log odds per standard deviation), the first region that applies is: (1) opposite sign, if the 95\% upper limit is below 0; (2) smaller than $\delta$, if the 90\% interval lies inside $(-\delta, \delta)$ (two one-sided permutation tests at $\pm\delta$); (3) holds, if the 95\% lower limit is above 0 and $\hat\Delta \geq \delta$; (4) attenuated, if the 95\% lower limit is above 0 and $\hat\Delta < \delta$; (5) inconclusive otherwise. A trained representation whose positive control failed was not interpretable. With 24 items there is no margin: opposite sign (upper limit $< 0$), positive (lower limit $> 0$), or inconclusive. The positive control required a negative Wald 95\% interval for the distance slope in the 8-item model of Section~\ref{sec:adults}. Claims that hold for all representations are intersection--union tests at each representation's own 95\% level; claims that the result depends on the representation are Holm-adjusted within two families per margin (trained representations other than DINOv2-small, two tests each; the two controls, two tests each). These $p$ values were never used to choose models.

\paragraph{Geometry calibration.} For each pair of representations, Pearson's $r$ was computed between their standardized distances over the 57{,}600 rows of the 8-item data, after residualizing both on level. Representations were clustered by average linkage (UPGMA) on $1 - r$ and the tree was cut at heights .3, .5, and .7. At a given cut the grouping \emph{discriminates} only if DINOv2-base falls in the group of DINOv2-small and neither control (random ResNet-50, pixels) shares a group with a trained encoder; otherwise that cut has ``no groups'' and no claim based on groups is made. DINOv2-base grouped with DINOv2-small at all three cuts, but a control shared a group with a trained encoder at every cut (MAE with pixels at .3; both controls with CVCL and MAE at .5 and .7), so calibration failed at every cut.

\paragraph{Regions for the CIFAR analyses (H1, H2, D).} With SESOI $\in \{.10, .20, .30\}$ on $|\rho|$ or the partial $|r|$, and without multiplicity adjustment, the first region that applies is: support, if $\rho > 0$ with Mantel $p < .05$ (H1), or the distance coefficient is negative with $p < .05$ (HC3 standard errors) and the semantic-distance coefficient has $p \geq .10$ or a partial $R^2$ below one third of the distance's (H2 and D); no support, if the 95\% interval lies entirely below the SESOI (H1), or excludes every value $\leq -\mathrm{SESOI}$ or shows an opposite effect with $p < .05$ (H2 and D); inconclusive otherwise. Support therefore does not depend on the SESOI, which affects only the boundary between no support and inconclusive. The plan anticipated that, with the half-width of the H2 interval near .35 on ImageNet-HSJ, no support would be nearly unattainable. Intervals for $r$ use the Fisher $z$ transformation with standard error $1/\sqrt{n - 4}$.

\paragraph{How $p < .05$ is used.} In the first plan, the six confirmatory tests (H1 on three datasets, H2 on two, and H4) were Holm-adjusted as one family; H5 and H6 were exploratory, with Benjamini--Hochberg correction. In the multi-encoder plans, per-representation regions use unadjusted $p < .05$ (or interval limits), and Holm adjustment applies only to the union claims described above. Observation-level likelihood-ratio tests, such as those of the trial-order models in Section~\ref{sec:adults}, are exploratory and do not enter any decision.
